\section*{Bab \ref{ch:b2:reproductive-hormones} --- Kendali Hormon atas Perkembangbiakan}

\begin{solution}{exo:b2:reproductive-hormones:1}
Hipotalamus (GnRH, dalam denyut) $\to$ hipofisis anterior (LH, FSH)
$\to$ testis: LH pada \omterm{def:b2:mammal-reproduction:testis}{sel Leydig} (testosteron), FSH pada \omterm{def:b2:mammal-reproduction:testis}{sel Sertoli}
(\omterm{def:b2:mammal-reproduction:testis}{spermatogenesis}, inhibin). Umpan baliknya: testosteron menghambat GnRH
dan LH; inhibin menghambat FSH.
\end{solution}

\begin{solution}{exo:b2:reproductive-hormones:2}
Ovariumnya: \omterm{def:b2:reproductive-hormones:cycle}{fase folikel}, hari 1--14, estradiol dari \omterm{def:b2:mammal-reproduction:ovary}{folikel} yang
tumbuh; \omterm{def:b2:mammal-reproduction:ovary}{ovulasi}, hari ke-14, \omterm{def:b2:reproductive-hormones:cycle}{lonjakan LH}; \omterm{def:b2:reproductive-hormones:cycle}{fase lutein}, hari 15--28,
progesteron dari \omterm{def:b2:mammal-reproduction:ovary}{korpus luteumnya}. Rahimnya: haid, hari 1--5
(penarikan steroidnya); fase proliferasi, hari 6--14 (estradiol); fase
sekresi, hari 15--28 (progesteron).
\end{solution}

\begin{solution}{exo:b2:reproductive-hormones:3}
Estrogen dan progestin hariannya menahan LH dan FSH tetap rendah lewat
umpan balik negatif, sehingga tidak ada \omterm{def:b2:mammal-reproduction:ovary}{folikel} yang direkrut, tidak
ada kenaikan estradiol yang terjadi, dan lonjakannya tidak pernah
menyala; progestinnya juga mengentalkan lendirnya. Pada pekan tanpa
pil, steroidnya turun dan \omterm{def:b2:reproductive-hormones:cycle}{endometriumnya}, yang dibangun di bawahnya,
luruh --- sebuah perdarahan penarikan, bukan sebuah haid sesudah
\omterm{def:b2:mammal-reproduction:ovary}{ovulasi}.
\end{solution}

\begin{solution}{exo:b2:reproductive-hormones:4}
LH dan FSH naik beberapa kali lipat (tanpa umpan balik testosteron
maupun inhibin); otot, gairah, dan kelenjar tambahannya menyusut.
Testosteron memulihkan cirinya dan membawa LH-nya turun kembali, tetapi
tidak kesuburannya: testisnya sudah tiada, dan bahkan pada seorang pria
yang masih bertestis, testosteron dari luar menekan LH dan dengan
demikian menekan testosteron di dalam testisnya yang dibutuhkan
\omterm{def:b2:mammal-reproduction:testis}{spermatogenesis}.
\end{solution}

\begin{solution}{exo:b2:reproductive-hormones:5}
$24/1.5 = 16$ denyut sehari; $24/4 = 6$ denyut. Ketika \omterm{def:b2:mammal-reproduction:ovary}{korpus luteumnya}
mati, denyutnya masih lambat sehingga memihak FSH: FSH naik lebih
dahulu --- hormon yang tepat, karena FSH-lah yang merekrut kelompok
\omterm{def:b2:mammal-reproduction:ovary}{folikel} berikutnya.
\end{solution}

\begin{solution}{exo:b2:reproductive-hormones:6}
\qty{200}{pg/mL} $= \qty{200}{ng/L}$: $200\times 10^{-9}/272 =
\qty{7.4e-10}{mol/L} = \qty{740}{pmol/L}$. \qty{15}{ng/mL} $=
\qty{15}{\micro g/L}$: $15\times 10^{-6}/314 = \qty{48}{nmol/L}$. Di
dalam satu mikroliter: $7.4\times 10^{-16}$ mol $= 4.4\times 10^{8}$
molekul.
\end{solution}

\begin{solution}{exo:b2:reproductive-hormones:7}
\omterm{def:b2:mammal-reproduction:ovary}{Ovulasi} pada hari $35 - 14 = 21$; subur dari hari ke-18 sampai hari
ke-22. Dengan daur yang tak teratur, \omterm{def:b2:reproductive-hormones:cycle}{fase folikelnya}, dan dengan
demikian hari \omterm{def:b2:mammal-reproduction:ovary}{ovulasinya}, tidak dapat diramalkan dari daur yang lalu,
dan jendelanya terlewat atau salah tempat.
\end{solution}

\begin{solution}{exo:b2:reproductive-hormones:8}
Dari hari ke-21 sampai ke-26 adalah 2,5 kali pelipatduaan:
$2^{2.5} = 5.7$ IU/L --- cukup. Berimplantasi pada hari ke-25 memberi
$1.4$ IU/L pada hari ke-26; \omterm{def:b2:mammal-reproduction:ovary}{korpus luteumnya}, yang sudah menyusut,
tidak terselamatkan, progesteronnya turun, dan lembaganya hilang
bersama haidnya. \omterm{def:b2:mammal-reproduction:implantation}{Implantasi} yang terlambat adalah penyebab lazim
kehilangan dini.
\end{solution}

\begin{solution}{exo:b2:reproductive-hormones:9}
Hipofisisnya menanggapi pola GnRH-nya, bukan jumlahnya: denyut tiap jam
mempertahankan LH dan FSH, sedangkan takaran harian yang sama yang
diinfuskan secara sinambung meniadakannya. Hal itu menyingkirkan
tanggapan takaran yang sederhana. Sebulan agonis yang bekerja lama:
sesudah nyala beberapa hari, LH-nya jatuh mendekati nol (\omterm{prop:b2:reproductive-hormones:pulses}{penyahpekaan
reseptor}), estradiolnya ke kadar pascamenopause, \omterm{def:b2:reproductive-hormones:cycle}{endometriumnya}
mengerut, dan tidak ada daur yang terjadi --- sebuah menopause medis
yang dapat balik.
\end{solution}

\begin{solution}{exo:b2:reproductive-hormones:10}
Tanpa hormon: $R^* = 100/0.5 = 200$. Sinambung: $100/6.5 = 15$. Sebuah
denyut dua menit membuang $6\times 200\times(2/60) = 40$ reseptor
(\qty{20}{\%}); dalam satu jam berikutnya kekurangannya mengerut dengan
faktor $e^{-0.5}$, yakni \qty{39}{\%} darinya dipulihkan. Lumbungnya
menetap di tempat kehilangan tiap denyutnya sama dengan pemulihan tiap
jamnya: kira-kira 150 reseptor sebelum tiap denyutnya, yaitu tiga
perempat maksimumnya --- sepuluh kali kadar sinambungnya.
\end{solution}

\begin{solution}{exo:b2:reproductive-hormones:11}
Menopause: tidak ada \omterm{def:b2:mammal-reproduction:ovary}{folikel}, sehingga tidak ada estradiol maupun
inhibin, sehingga tidak ada umpan balik: LH dan FSH-nya berjalan
tinggi. Prolaktinoma: prolaktinnya memperlambat pembangkit denyut
GnRH-nya, sehingga LH-nya rendah, \omterm{def:b2:mammal-reproduction:ovary}{folikelnya} tidak digerakkan,
estradiolnya rendah, dan tidak ada daur. Tumor sel granulosa: estradiol
tinggi yang tetap menekan LH lewat umpan balik negatif (dan tanpa
\omterm{def:b2:mammal-reproduction:ovary}{folikel} yang matang tidak ada yang dapat diovulasikan), sehingga
daurnya berhenti dengan estradiol tinggi --- dan \omterm{def:b2:reproductive-hormones:cycle}{endometriumnya} tumbuh
berlebihan.
\end{solution}

\begin{solution}{exo:b2:reproductive-hormones:12}
Estradiol yang rendah atau singkat menurunkan GnRH dan LH; estradiol
yang tinggi dan bertahan (di atas \qty{200}{pg/mL} selama lebih dari 36
jam) mengubah tanggapan sistem kisspeptin--GnRH dan tanggapan
hipofisisnya, yang sementara itu telah menimbun LH, sehingga hormon
yang sama kini membangkitkan sebuah ledakan. \omterm{def:b2:mammal-reproduction:ovary}{Ovulasi} haruslah
serba-atau-tidak dan tepat waktunya: sel telurnya harus dilepaskan
sekali, pada satu hari, ketika sebuah \omterm{def:b2:mammal-reproduction:ovary}{folikel} masak. Sebuah tombol
putar (LH sebanding dengan estradiol) akan memberi peluteinan yang
berangsur dan sebagian, \omterm{def:b2:mammal-reproduction:ovary}{folikel} yang berovulasi setengah masak atau
beberapa sekaligus, dan tidak ada hari subur yang jelas.
\end{solution}

\begin{solution}{pb:b2:reproductive-hormones:1}
\textbf{1.} \omterm{def:b2:reproductive-hormones:cycle}{Fase folikel}, hari 1--14: estradiolnya naik, progesteronnya
rendah. \omterm{def:b2:reproductive-hormones:cycle}{Fase lutein}, hari 15--28: progesteronnya tinggi.
\textbf{2.} Lonjakannya memuncak pada hari ke-14 (mulai pada hari
ke-13); \omterm{def:b2:mammal-reproduction:ovary}{ovulasinya} kira-kira 36 jam sesudah awalnya, pada hari ke-15.
\textbf{3.} Hari ke-15 sampai ke-28: 14 hari, yaitu umur \omterm{def:b2:mammal-reproduction:ovary}{korpus
luteumnya}.
\textbf{4.} Progesteron menaikkan titik acuan pengaturan suhunya
\qty{0.3}{\celsius} sampai \qty{0.4}{\celsius}; kenaikannya muncul satu
atau dua hari sesudah \omterm{def:b2:mammal-reproduction:ovary}{ovulasi}, sehingga ia menegaskan bahwa \omterm{def:b2:mammal-reproduction:ovary}{ovulasinya}
telah terjadi tetapi tidak dapat mengumumkannya.
\textbf{5.} Hari ke-11 sampai ke-14 (210, 260, 310, 220): tiga sampai
empat hari, lebih daripada 36 jam yang dituntut sakelarnya.
\textbf{6.} Hari ke-12 sampai ke-16.
\textbf{7.} $310\times 10^{-9}/272 = \qty{1.14e-9}{mol/L} =
\qty{1140}{pmol/L}$; $16\times 10^{-6}/314 = \qty{51}{nmol/L}$.
\textbf{8.} $65/8 = 8$. Dengan waktu paruh satu jam, LH-nya memeroh
tiap jam begitu sekresinya berhenti, sehingga sehari kemudian ia
kembali mendekati garis dasarnya.
\textbf{9.} \omterm{def:b2:mammal-reproduction:ovary}{Ovulasi} sekitar hari ke-22, daur kira-kira 36 hari.
\textbf{10.} \omterm{def:b2:mammal-reproduction:implantation}{Implantasi} pada hari ke-22; pada hari ke-27 (lima hari,
2,5 kali pelipatduaan) kira-kira \qty{5.7}{IU/L}.
\textbf{11.} Progesteron yang turun ke \qty{2}{ng/mL} pada hari ke-27
berarti \omterm{def:b2:mammal-reproduction:ovary}{korpus luteumnya} menyusut sesuai jadwal: tidak ada yang
menyelamatkannya.
\textbf{12.} Progesteron sekitar \qty{20}{ng/mL} dan naik, estradiol
sekitar \qty{200}{pg/mL}.
\textbf{13.} Pada akhir daurnya, estradiol, progesteron, dan inhibinnya
turun bersama-sama dan FSH-nya lolos dari umpan baliknya; sembari
kelompoknya tumbuh, estradiol dan inhibinnya naik dan FSH-nya turun,
dan hanya \omterm{def:b2:mammal-reproduction:ovary}{folikel} dengan reseptor FSH terbanyak (sel granulosa
terbanyak) yang terus tumbuh pada FSH yang menurun --- sisanya mati.
\textbf{14.} $R^* = 200$ tanpa hormon; $100/6.5 = 15$ di bawah hormon
sinambung.
\textbf{15.} $6\times 200\times 2/60 = 40$ reseptor, yaitu
\qty{20}{\%}; tetapan waktunya $1/d = \qty{2}{h}$, sehingga
\qty{39}{\%} kekurangannya dipulihkan dalam satu jam; lumbungnya
menetap sekitar 150 reseptor, yaitu tiga perempat maksimumnya.
\textbf{16.} Denyutnya menjaga hipofisisnya pada kira-kira 150 reseptor
dan ia menjawab tiap denyutnya dengan LH; sebuah infus sinambung
mendorong lumbungnya ke 15 dan sel yang sama membisu, berapa pun
takarannya; denyutnya memulihkan lumbung dan tanggapannya.
\textbf{17.} Hari pertama: LH dan testosteronnya naik (agonisnya
merangsang reseptor yang masih ada --- itulah nyalanya). Satu bulan:
reseptornya ternyahpekakan, LH-nya mendekati nol, testosteronnya pada
kadar pengebirian, dan itulah tujuan pengobatannya.
\textbf{18.} FSH-nya naik (tanpa inhibin), LH-nya tetap normal; lebih
banyak \omterm{def:b2:mammal-reproduction:ovary}{folikel} yang selamat dari pemilihan tiap daurnya, dengan \omterm{def:b2:mammal-reproduction:ovary}{ovulasi}
berganda dan kembar dizigot.
\textbf{19.} Pada daur alaminya, penurunan FSH membunuh semua kecuali
folikel dominannya; pasokan FSH yang tetap menjaga seluruh kelompoknya
tumbuh sampai matang.
\textbf{20.} Sepuluh \omterm{def:b2:mammal-reproduction:ovary}{folikel} membuat jauh lebih banyak daripada
\qty{200}{pg/mL} estradiol secara dini, dan itu akan memicu \omterm{def:b2:reproductive-hormones:cycle}{lonjakan LH}
yang terlalu dini serta \omterm{def:b2:mammal-reproduction:ovary}{ovulasi} sebelum pengumpulannya; antagonisnya
menghadang GnRH dan tidak ada lonjakan yang dapat terjadi.
\textbf{21.} Pada tabelnya, lonjakannya memuncak pada hari ke-14 dan
\omterm{def:b2:mammal-reproduction:ovary}{ovulasinya} menyusul kira-kira 36 jam sesudah awalnya; hCG menggantikan
lonjakannya, dan pengumpulan pada 35 jam menangkap oositnya dalam
keadaan matang tetapi masih di dalam \omterm{def:b2:mammal-reproduction:ovary}{folikelnya}.
\textbf{22.} $10\times 0.7 = 7$ lembaga, $7\times 0.4 = 2.8$
\omterm{def:b2:mammal-reproduction:implantation}{blastosista}; $1 - 0.65^{2} = 0.58$.
\textbf{23.} Estrogen dan progestin yang tetap menjaga FSH dan LH tetap
rendah: tidak ada \omterm{def:b2:mammal-reproduction:ovary}{folikel} yang tumbuh, estradiolnya tidak pernah naik
sampai ambangnya, tidak ada lonjakan, tidak ada \omterm{def:b2:mammal-reproduction:ovary}{korpus luteum}, dan
tidak ada kenaikan progesteron.
\textbf{24.} Testosteron harus diberikan bagi ciri sekunder dan
gairahnya; tetapi ia tidak memulihkan kepekatan tinggi di dalam
testisnya yang disediakan \omterm{def:b2:mammal-reproduction:testis}{sel Leydig}, sehingga \omterm{def:b2:mammal-reproduction:testis}{spermatogenesisnya} tetap
tertekan --- dan itulah maksudnya.
\textbf{25.} \omterm{def:b2:mammal-reproduction:ovary}{Ovulasi} pada hari ke-15; \omterm{def:b2:reproductive-hormones:cycle}{fase lutein} 14 hari; jendela
subur hari ke-12 sampai ke-16; lonjakannya dipicu estradiol di atas
\qty{200}{pg/mL} selama lebih dari 36 jam, yang terpenuhi pada hari
ke-11 sampai ke-14.
\end{solution}
